\documentclass[10pt]{beamer}

\usetheme{metropolis}
\usepackage{appendixnumberbeamer}

\usepackage{booktabs}
\usepackage[scale=2]{ccicons}

\usepackage{pgfplots}
\usepgfplotslibrary{dateplot}

\usepackage{xspace}
\newcommand{\themename}{\textbf{\textsc{metropolis}}\xspace}

\usepackage{listings}
\hypersetup{
    urlcolor=cyan,
}

\title{Git good}
\subtitle{Git tips and workflow}
\date{\today}
\author{Viktor Holmgren}
\institute{LiTHe kod}
\titlegraphic{\hfill\includegraphics[height=1.5cm]{lithekod-logo.pdf}}

\begin{document}

\maketitle

\begin{frame}{Table of contents}
  \setbeamertemplate{section in toc}[sections numbered]
  \tableofcontents[hideallsubsections]
\end{frame}

\section{Match that flow}

\begin{frame}[fragile]{Master and develop}
  Git Flow workflow typically recommends \emph{two main branches};
  \begin{description}
    \item master: official releases
    \item develop: integration branch for features
  \end{description}

  \begin{lstlisting}
  git branch develop
  git push -u origin develop
  \end{lstlisting}

  \includegraphics[width=\textwidth]{developbranch}

\end{frame}
\begin{frame}[fragile]{Use them branches yo}
  Using Topic branches (feature/bugfix) is an extremely good practice!
  \begin{itemize}
    \item Create a new branch from develop for each new \emph{feature}
    \item Develop your feature
    \item Interactive rebase if needed to clean up code
    \item Merge into develop using CI
  \end{itemize}

  \includegraphics[width=\textwidth]{feature}
\end{frame}

\begin{frame}[fragile]{Lets ship it!}

  \begin{itemize}
    \item When develop has enough features, or release is apporaching
    \item Fork new branch from develop
    \item Continue with bugfixes, documentation but \textbf{no features} 
    \item When ready merge \emph{release} into \emph{master} and tag, merge \emph{master} into \emph{develop}
  \end{itemize}

  \textbf{Good}: Work on release and develop in parallel, easy to see what realease is in process

  \includegraphics[width=0.75\textwidth]{release}
\end{frame}

\section{Making it look nice af}
\begin{frame}{Patch it up baby}
  \begin{columns}[T] % contents are top vertically aligned
     \begin{column}[T]{.5\textwidth} % each column can also be its own environment
      Git stores changes in three separate states:

      \begin{itemize}
        \item {Working directory}: Your files locally, currently.
        \item {Staged/Index}: Changes to be committed.
        \item {Committed}: Changes stored in the history.
      \end{itemize}

      \includegraphics[width=\textwidth]{status}

     \end{column}
     \begin{column}[T]{.5\textwidth} % alternative top-align that's better for graphics
      Use \textit{git status} and git will tell you what to do!
      \begin{itemize}
        \item \textit{git add} add files to index, \textit{git add -p} adds \textit{parts} of a file
        \item \textit{git reset} removes files from index, \textit{git reset -p} removes \textit{parts} of a file
        \item \textit{git reset --soft} removes commits and moves changes to index 
        \item \textit{git reset --hard} removes commits and discards changes.
      \end{itemize}
     \end{column}
   \end{columns}
\end{frame}

\begin{frame}{Stashing them goodies}
  Sometimes you need to switch to work on something else, or store an attempt.
  Either use a \textit{Work in progress} branch or stash your changes.
  \begin{itemize}
  \item WIP: \textit{git checkout -b WIP-<something> \&\& git add . \&\& git commit -m 'WIP'}
    \item Stash: \textit{git stash} then later \textit{git stash pop}
  \end{itemize}
\end{frame}

\begin{frame}[fragile]{What the rebase?}
  \textit{Rebasing} is an alternative to \textit{merging}.

  Will move commits from a given branch to the current branch.

  \begin{lstlisting}
  git rebase <branch>
  \end{lstlisting}

  \begin{itemize}
		\item Results in a linear history
		\item Will rewrite history (different history results in different commits)
	\end{itemize}

  \textbf{You shallt not rebase any shared changes}

  With great power comes great responsibility
  \includegraphics[width=0.75\textwidth]{rebase}
\end{frame}

\begin{frame}[fragile]{What kind of sorcery is this?}
  Select and transform commits arbitrarily 

  \begin{lstlisting}
  git rebase -i <branch>
  \end{lstlisting}

  \begin{itemize}
		\item Remove or add commits at any point
		\item Change order of commits
    \item Squash commits together
		\item Edit commit messages at any point
	\end{itemize}
\end{frame}

\begin{frame}[plain]{What kind of sorcery is this?}
  \includegraphics[width=\textwidth]{rebasing}
\end{frame}

\begin{frame}{But why?}
  \textbf{Please use before you feature/bugfix branch to remote}.
  \begin{itemize}
    \item Removing ulgy: fix <insert minor problem> commits by squash
    \item Undoing part of the work
    \item Splitting one large commit into smaller parts
    \item Combining smaller relevant parts
		\item Fix spelling mistakes in commit messages
    \item Correcing styleguide or adding comments in the same commit
    \item ...
	\end{itemize}
\end{frame}

\section{What the hell is happening?!}
\begin{frame}[fragile]{Logging like a bawse}
  \begin{lstlisting}
    git log --decorate --oneline --graph
  \end{lstlisting}

  \begin{columns}[T] % contents are top vertically aligned
     \begin{column}[T]{.5\textwidth} % each column can also be its own environment
       \includegraphics[width=\textwidth]{bad-log}
       Bad log
     \end{column}
     \begin{column}[T]{.5\textwidth} % alternative top-align that's better for graphics
       \includegraphics[width=\textwidth]{good-log}
       Good log
     \end{column}
   \end{columns}
\end{frame}

\begin{frame}[fragile]{Show me your secrets}

  Show the changes made in a given commit:
  \begin{lstlisting}
    git diff <commit>^!
  \end{lstlisting}

  Show the staged changes (what will be committed)
  \begin{lstlisting}
    git diff --staged
  \end{lstlisting}
\end{frame}

\begin{frame}[fragile]{Blame it on the new guy}

  Git blame, aka 'Vem var det som kasta!'\\
  Show the changes made in a given commit:
  \begin{lstlisting}
    git blame <file>
  \end{lstlisting}

  \includegraphics[width=\textwidth]{blame}
\end{frame}

\section{Welp that was stupid}
\begin{frame}{Do it right, right a way}
  \begin{itemize}
    \item \textbf{Use an ignorefile!}, you probably want to set it up before anything.\\
      \href{https://gitignore.io/}{Gitignore.io} is really good.
    \item A global ignorefile can be good also if you work in similar languages or types of projects.\\
      Simply create an ignore file at $~/.gitignore\_global$\\
      Then run $git config --global core.excludesfile ~/.gitignore\_global$
  \end{itemize}
\end{frame}

\section{Bisecting FTW}
\begin{frame}[fragile]{Finding that pesky bug}
  Git bisect is a general command to binary search for \textbf{any} property, but most commonly a bug.

  \begin{lstlisting}
> git bisect start
> git bisect bad
> git bisect good e7f9b07
Bisecting: 13 revisions left to test
after this (roughly 4 steps)
[ea7de...] Implement event forwarding
  \end{lstlisting}
  Then run \textit{git bisect good} if that version works, or \textit{git bisect bad} if it does not work until you are left with the commit you were searching for.
\end{frame}

\section{Teamwork makes the dream work}

\begin{frame}[fragile]{Try to make sense, okay}
  \begin{columns}[T] % contents are top vertically aligned
     \begin{column}[T]{.5\textwidth} % each column can also be its own environment
      Writing a good commit message is art.
      But here are some guidelines:

      \begin{itemize}
        \item Limit the subject line to 50 characters
        \item Capitalize the subject line
        \item Do not end the subject line with a period
        \item Use the imperative mood in the subject line\\
          'Refactor A by moving Y' not 'Refactored A by moving Y'\\
          I.e like a command!
      \end{itemize}

      \url{https://chris.beams.io/posts/git-commit/}
     \end{column}
     \begin{column}[T]{.5\textwidth} % alternative top-align that's better for graphics
      \includegraphics[width=\textwidth]{commit-logs}

     \end{column}
   \end{columns}

\end{frame}

\begin{frame}[fragile]{Code reviews, code review, code reviews}
  \begin{columns}[T] % contents are top vertically aligned
     \begin{column}[T]{.5\textwidth} % each column can also be its own environment
      Use code reviews for everything, and I mean \textbf{everything}!\\
      We used it to great effect in writing reports!\\
      For each topic branch, use some type of code review.

      You will:
      \begin{itemize}
        \item Spread knowledge about what is being worked on, and changed (knowledge sharing)
        \item Catch bugs and faults early!
      \end{itemize}

     \end{column}
     \begin{column}[T]{.5\textwidth} % alternative top-align that's better for graphics
      \includegraphics[width=\textwidth]{pum-example}
     \end{column}
   \end{columns}

\end{frame}

\section{Questions and demo?}

\end{document}
